% Section 6: Type Exclusion and C*-Algebra Promotion
% \input-able LaTeX file for Paper 5
% Conventions: a \sp b = sequential product (dot notation matching vdW)
%   JVW = Jordan-von Neumann-Wigner classification
%   X* = X^dagger = conjugate transpose

\section{Type Exclusion and $C^*$-Algebra Promotion}\label{sec:exclusion}

%--------------------------------------------------------------------
\subsection{The Jordan--von Neumann--Wigner
Classification}\label{subsec:jvw}

By Theorem~\ref{thm:EJA}, the state space~$V$ is a Euclidean Jordan
algebra (EJA). The Jordan--von Neumann--Wigner
classification~\cite{JordanVonNeumannWigner1934} asserts that every finite-dimensional
simple EJA belongs to one of five families:

\begin{center}
\begin{tabular}{lllcl}
\toprule
Type & Notation & $\dim$ & LT? & Exclusion mechanism \\
\midrule
Real & $M_n(\mathbb{R})^{\mathrm{sa}}$ & $n(n{+}1)/2$
  & Only $n{=}1$ & Dimension mismatch \\
Complex & $M_n(\mathbb{C})^{\mathrm{sa}}$ & $n^2$
  & All $n$ & --- \\
Quaternionic & $M_n(\mathbb{H})^{\mathrm{sa}}$ & $n(2n{-}1)$
  & Only $n{=}1$ & Dimension mismatch \\
Spin factor & $V_n$ \,($n \ge 3$) & $n{+}1$
  & Only $V_3$ & Barnum--Wilce \\
Albert & $M_3(\mathbb{O})^{\mathrm{sa}}$ & $27$
  & Never & No composite exists \\
\bottomrule
\end{tabular}
\end{center}

\noindent
The spin factors overlap with the matrix types at low dimensions:
$V_2 \cong M_2(\mathbb{R})^{\mathrm{sa}}$,
$V_3 \cong M_2(\mathbb{C})^{\mathrm{sa}}$, and
$V_5 \cong M_2(\mathbb{H})^{\mathrm{sa}}$.

%--------------------------------------------------------------------
\subsection{Excluding Non-Complex Types}\label{subsec:exclusion}

We now show that local tomography (Theorem~\ref{thm:local-tomo})
excludes every non-complex type.

\medskip
\noindent
\textbf{Real matrices, $M_n(\mathbb{R})^{\mathrm{sa}}$ ($n \ge 2$).}
The local tomography condition $[\dim(V)]^2 = \dim(\text{composite})$
becomes
\[
  \left[\tfrac{n(n+1)}{2}\right]^2
  = \tfrac{n^2(n^2+1)}{2}\,.
\]
Simplifying: $(n+1)^2 = 2(n^2+1)$, which reduces to
$(n-1)^2 = 0$, forcing $n = 1$ (trivial,
one-dimensional). For $n \ge 2$, the real composite is strictly larger
than the product-effect span: $d^2 < \dim(\text{composite})$, and product
measurements leave an entangled sector
unresolved~\cite{BarnumWilce2014}.

\medskip
\noindent
\textbf{Quaternionic matrices,
$M_n(\mathbb{H})^{\mathrm{sa}}$ ($n \ge 2$).}
The analogous condition yields
$[n(2n-1)]^2 = n^2(2n^2-1)$, which again reduces to
$(n-1)^2 = 0$, forcing $n = 1$ (trivial). For $n \ge 2$,
the quaternionic composite has \emph{fewer} dimensions than $d^2$:
product effects are linearly dependent on the joint
system~\cite{BarnumWilce2014}.

\medskip
\noindent
\textbf{Spin factors, $V_n$ ($n \ge 4$).}
The cross-identifications $V_2\cong M_2(\mathbb{R})^{\mathrm{sa}}$ and
$V_5\cong M_2(\mathbb{H})^{\mathrm{sa}}$ are already excluded above.
For $V_n$ with $n\ge 4$ and $n\ne 5$, these are ``pure'' spin factors not
isomorphic to any $M_k(\mathbb{K})^{\mathrm{sa}}$.
Barnum and Wilce~\cite{BarnumWilce2014} prove that no locally tomographic
composite exists for these types (Theorem~4.1 of~\cite{BarnumWilce2014}).

\medskip
\noindent
\textbf{Albert algebra, $M_3(\mathbb{O})^{\mathrm{sa}}$.}
The exclusion here is the strongest: Barnum, Graydon, and
Wilce~\cite{BarnumUdudeckVdW2020} prove that no non-signaling composite with product
states exists for any EJA containing an Albert algebra summand
(Theorem~1 of~\cite{BarnumUdudeckVdW2020}). The 27-dimensional exceptional Jordan
algebra cannot form \emph{any} consistent bipartite system, even a
non-locally-tomographic one.

\medskip
\noindent
\textbf{Complex matrices, $M_n(\mathbb{C})^{\mathrm{sa}}$ ($n \ge 1$).}
$\dim(M_n(\mathbb{C})^{\mathrm{sa}}) = n^2$, and the composite
$M_{n^2}(\mathbb{C})^{\mathrm{sa}}$ has dimension $(n^2)^2 = n^4 =
(n^2)^2$. Local tomography holds for all~$n$: the dimension formula
$d = n^2$ is \emph{multiplicatively stable}. The complex type is the
unique survivor.

%--------------------------------------------------------------------
\subsection{C*-Algebra Promotion}\label{subsec:cstar}

The surviving EJA is promoted to a C*-algebra by a single published
theorem.

\begin{theorem}[{C*-algebra structure;
  \cite[Theorem~3]{vandeWetering2019b}}]\label{thm:vdW3}
Let $V$ be a sequential product space satisfying S1--S7. If the
composite $V\otimes V$, equipped with the product-form sequential
product, is also a sequential product space and the composite is locally
tomographic, then $V$ is the self-adjoint part of a C*-algebra.
\end{theorem}

\begin{center}
\begin{tabular}{lll}
\toprule
\textbf{Hypothesis} & \textbf{Status} & \textbf{Established in} \\
\midrule
$V$ satisfies S1--S7
  & Verified & Section~\ref{sec:axioms} \\
$V\otimes V$ with product-form SP satisfies S1--S7
  & Verified & Proposition~\ref{prop:inheritance} \\
Composite is locally tomographic
  & Verified & Theorem~\ref{thm:local-tomo} \\
\bottomrule
\end{tabular}
\end{center}

\noindent
Theorem~\ref{thm:vdW3} combined with the type exclusion of
Section~\ref{subsec:exclusion} establishes that
$V \cong M_n(\mathbb{C})^{\mathrm{sa}}$ for some $n \ge 2$.  This is
the self-adjoint part of the C*-algebra $M_n(\mathbb{C})$, and the
C*-promotion is complete.

\begin{remark}[Consistency with independent results]\label{rem:consistency}
Two further published theorems provide independent consistency checks
on the conclusion $V \cong M_n(\mathbb{C})^{\mathrm{sa}}$.

\emph{Barnum--Wilce~\cite{BarnumWilce2014}.}
Barnum and Wilce prove that in a locally tomographic dagger-compact
probabilistic theory containing a qubit, every EJA system type is
$M_n(\mathbb{C})^{\mathrm{sa}}$.  Their hypotheses are categorical
(they require a full monoidal theory with dagger structure on every
object, not merely one composite), and we do not construct this
categorical apparatus here.  However, our conclusion is consistent
with their result: after establishing $V = M_n(\mathbb{C})^{\mathrm{sa}}$
via Theorem~\ref{thm:vdW3}, the category of complex matrix algebras
with Hilbert space tensor products does form a locally tomographic
dagger-compact theory.

\emph{Hanche-Olsen~\cite{HancheOlsen1985}.}
Hanche-Olsen proves that a JB-algebra $E$ whose tensor product
$E \otimes E_2$ with the qubit algebra
$E_2 = M_2(\mathbb{C})^{\mathrm{sa}}$ is again a JB-algebra
(with factorization identities for the Jordan product) must be the
self-adjoint part of a C*-algebra.  The qubit subsystem $E_2$ is
available only \emph{after} type exclusion has already established
the complex type, so this theorem does not provide an independent
route to the conclusion.  It does, however, confirm the C*-algebra
structure through entirely different mathematical machinery
(JB-algebra tensor products rather than sequential product
classification).

Neither theorem is used in the derivation; both serve as cross-checks
that the conclusion is robust.
\end{remark}

%--------------------------------------------------------------------
\subsection{The Involution}\label{subsec:involution}

Theorem~\ref{thm:vdW3} and the type exclusion of
Section~\ref{subsec:exclusion} establish that
$V = M_n(\mathbb{C})^{\mathrm{sa}}$ is the self-adjoint part of the
C*-algebra $A = M_n(\mathbb{C})$. The C*-algebra involution is the
conjugate transpose:
\[
  X^* \;=\; X^\dagger \;=\; \overline{X}^{\,T}.
\]
This involution satisfies four defining properties:
\begin{enumerate}
\item[\emph{(P1)}] \textbf{Involutive:} $(X^*)^* = X$.
\item[\emph{(P2)}] \textbf{Anti-multiplicative:} $(XY)^* = Y^*X^*$.
\item[\emph{(P3)}] \textbf{Fixed-point set:}
  $\{X \in M_n(\mathbb{C}) : X^* = X\}
  = M_n(\mathbb{C})^{\mathrm{sa}}$.
\item[\emph{(P4)}] \textbf{C*-identity:}
  $\|X^*X\| = \|X\|^2$ (operator norm).
\end{enumerate}
Property (P3) identifies the self-adjoint part recovered from the
self-modeling construction with the standard notion of Hermitian matrices.
Property (P4) equips $M_n(\mathbb{C})$ with its C*-norm, completing the
C*-algebra structure.

%--------------------------------------------------------------------
\subsection{Main Theorem}\label{subsec:main-theorem}

We can now state the central result.

\begin{theorem}[Self-modeling implies complex quantum
mechanics]\label{thm:main}
Let $(V, \le, \mathbf{1})$ be a finite-dimensional spectral order unit
space. If $V$ admits a faithful self-model $\varphi\colon V \to V_M$
with $V_M$ an isomorphic copy, and if the minimal composite $V_{BM}$
satisfies non-signaling and product-form sequential product inheritance,
then:
\begin{enumerate}
\item $V \cong M_n(\mathbb{C})^{\mathrm{sa}}$ for some $n \ge 2$;
\item the sequential product is the L\"uders product
  $\sp{a}{b} = \sqrt{a}\, b\, \sqrt{a}$;
\item $M_n(\mathbb{C})$ carries the conjugate-transpose involution
  $X^* = X^\dagger$.
\end{enumerate}
\end{theorem}

In particular, the complex field, the Jordan product, the C*-involution,
and local tomography are all \emph{derived} from the single operational
premise of faithful self-modeling.

\begin{proof}[Proof outline]
\mbox{}%
\begin{enumerate}
\item The self-modeling constraint yields a sequential product on
  $V$ via the ``test--update--test'' operational
  semantics (Section~\ref{sec:sp}).
\item The corrected product formula~\eqref{eq:corrected-product},
  incorporating the Peirce $1$-space feedback selected by faithful
  tracking, satisfies S1--S7
  (Section~\ref{sec:axioms}).
\item By Theorem~\ref{thm:EJA} (van de Wetering), $V$ is an EJA.
\item The composite $V_{BM}$ with the product-form SP inherits
  S1--S7 (Proposition~\ref{prop:inheritance}), and state
  separation combined with minimality yields local tomography
  (Theorem~\ref{thm:local-tomo}).
\item Dimension counting excludes all non-complex simple EJA types
  (Section~\ref{subsec:exclusion}).
\item Theorem~\ref{thm:vdW3} (C*-promotion via van de Wetering)
  yields $V = M_n(\mathbb{C})^{\mathrm{sa}}$ with the
  conjugate-transpose involution. \qed
\end{enumerate}
\end{proof}
